\section{Experimental Method}
\label{sec:evaluation}
The study asks three questions: how much of a burst reduction comes from work
placement within a fixed implementation; how that choice changes with native
kernels and completion horizon; and whether it improves complete-module
behavior under Rack's sample-synchronous engine. These questions and the
workload matrix were recorded before collection. The present analysis is a
\emph{descriptive pilot study}; it does not select and then confirm a winner.

\subsection{Host, sessions, and workload boundaries}
Two separately prepared sessions ran on September 30, 2026 (local time) on one
Apple M1 Pro laptop, with macOS 26.6.2 and Apple Clang 21.0.0. Both used the
same executable and archived sources. The operator declared networking and
Bluetooth disconnected and agents and unnecessary applications closed.
The runner required AC power and Low Power Mode off, held
\code{caffeinate -is} assertions, settled for 180 seconds after staging, and
settled each fresh process for one second before warmup. Pre/post snapshots
verify power and sleep assertions; they also show residual macOS service
activity, including \code{sandboxd} and Spotlight indexing. They cannot
establish the absence of interference during a particular interval.

Each session executed 194 planned cells twice in fresh processes, giving 776
processes in total. Three bridge cells repeat between groups: there are 191
unique full configurations. Group order was fixed; within-group process order
used the predeclared session seeds. We retain every process and slow interval.
Table~\ref{tab:study-design} describes the seven groups. Setup, plan creation,
correctness replay, storage probes, and output serialization are outside the
block timing intervals. These measurements replace the earlier manuscript's
performance campaign; old observations are retained only as historical artifacts.

\begin{table}[htbp]
\centering\small
\begin{tabular}{@{}lrll@{}}
\toprule
Group & Cells & Execution & Blocks/process \\
\midrule
Baselines & 24 & continuous & 4096 \\
Granularity & 39 & continuous & 960--4096 \\
Scaling & 16 & continuous & 1024--16384 \\
Modules & 9 & paced & 960--5760 \\
Stress & 12 & continuous & 1028--16384 \\
Host & 83 & paced & 1024--4096 \\
Lifecycle & 11 & continuous & 1024--16384 \\
\bottomrule
\end{tabular}
\caption{Predeclared pilot matrix: two sessions and two fresh processes per cell per session. Stream groups use 256 nominal hops, except Stress (128); host uses 262,144 engine samples. Lifecycle includes separate 65,536-sample event sequences and 1,048,576-sample decays. Each process has 16 warmup hops; host warmup rounds up to whole blocks. Three repeated bridge cells retain their group identities.}
\label{tab:study-design}
\end{table}


The stream harness measures complete analysis calls with a controlled sink,
or actual module \code{process()} paths including conditioning and display-data
conversion. Fourier's shipped defaults use four SIMD lanes, $N=2048$, a
30\,ms hop ($H=1440$ at 48\,kHz), and a flattop window. Spectre uses scalar
samples, $N=2048$, $H=1024$, and its flattop default. Controlled scalar
comparisons use a Hann window, $H=1024$, no smoothing, and $D=64$ unless
specified. All figures and tables state which boundary they compare.

The host group calls the pinned Rack engine's actual \code{stepBlock()},
including worker synchronization, with one or four engine threads, 0/1/4/16
analyzers, and 0/16/64 deterministic background modules. The host matrix is a
structured set of combinations, not a full factorial design. Native module
controls replace analysis inside the same module shell while retaining input
conditioning, display conversion, and mailbox publication. They retain some
unused production storage; their footprint is not an optimized replacement's
memory requirement. A concurrent headless consumer optionally copies and
checksums snapshots at 30 or 60\,Hz, with declared stalls. It does not draw a
window or drive an audio interface.

\subsection{Matched controls and native layouts}
The matched first-party batch and distributed controls use the same scalar
analysis implementation, arithmetic, and storage. They isolate placement more
closely than a comparison across FFT libraries. Native batch/hybrid pairs also
share the same retained-input, window, magnitude, smoothing, and output path.
Batch runs the sequence immediately; hybrid distributes its surrounding stages
but executes one indivisible native transform/reorder call. The native
pipeline uses contiguous ring segments and native output layouts, avoiding
per-element ring-modulo and canonical-complex conversion passes. vDSP uses
vector window/packing operations and batched four-channel transforms; FFTW
uses reusable real plans, including \code{plan\_many} for four channels;
PFFFT uses its ordered packed spectrum. Magnitudes use an explicitly scaled
norm to prevent avoidable squaring underflow. Thus these are stronger native
controls than the previous scalar-glue adapters, but not a claim to exhaust
every provider-specific optimization.

The granularity group holds frame hop fixed and changes native completion
horizon to $H_c\in\{H,\ceil{H/2},\ceil{H/4}\}$. Publication occurs at
endpoint age $H_c-1$, followed by idle capture until the next endpoint. This
changes latency and surrounding-stage density without subdividing the native
FFT. Native sub-FFT leaves and new stage weights were explicitly deferred;
neither has measured evidence in this study. Production weights $p=4$ for
window rebuild and $o=2$ for Fourier output were tuned in earlier engineering
experiments on this same M1 Pro. Band-cache rebuilding remains inside output
units without an additional weight. The new pilots are not independent
validation of those weight choices.

\subsection{Timing, replication, and age}
Continuous stream groups execute blocks back to back. The host group uses
absolute periodic releases at $D/f_s$, without resetting the release origin
after overruns. A monotonic clock records wake, start, and finish separately.
Inherited thread policy is used: no core affinity, frequency lock, or
real-time policy experiment is implied. Stream processes inherit and record
FPU control; the actual Rack engine applies Rack's block FPU policy. Empty
timer observations are retained without subtraction, and values are rounded
to approximately three significant figures.

The primary burst statistic is the median of \emph{per-hop peak block times}:
for each complete endpoint-to-next-endpoint interval, take the largest duration
of any whole block intersecting it. This includes a batch burst regardless of
its frequency among all callbacks. Shared blocks contribute to each intersected
hop; no sub-block costs are invented. Engine peaks use actual replay endpoints
and the full hop for both batch and delayed publication. Lifecycle sequences
are excluded from this stationary-hop statistic. Callback p99, maxima, and
counts above $D/f_s$ remain separate secondary observations.

Stream cost is the sum of instrumented block wall times divided by engine
samples. It includes timer/dispatch and interruption effects; this new design
contains no separate untimed-loop throughput pass. Host results additionally
retain process-wide user-plus-system CPU time across the measurement interval,
including workers, pacing overhead and any headless consumer. CPU and block
wall time are distinct quantities. A paced release miss means finish exceeds
the synthetic release deadline, including wake lateness and accumulated
backlog. A compute-budget exceedance means the timed block alone exceeds
$D/f_s$. Neither is an audio-device underrun.

For each cell we average the two process costs within each session; for burst
quantiles we take the median of the two process summaries. Tables show the
midpoint of the two session summaries, their range where indicated, and the
largest retained observation. There are only two session replicates, both on
the same date and host. Callbacks, hops, analyzers, and channels are dependent
observations. We report descriptive variation rather than confidence intervals,
rare-event probabilities or a worst-case execution-time claim
\cite{kalibera2013,mytkowicz2009,wilhelm2008}.

Publication age is measured in logical engine samples relative to the frame
endpoint; center age adds $(N-1)/2$. Consumer age is bounded by the completed
block watermarks observed around its copy, with one block of uncertainty.
Final drains are excluded from normal consumer ages. These sample-clock ages
are not wall-clock freshness during an overrun and do not include repaint.

\subsection{Numerical and provenance checks}
Every session passed checks of archive and artifact hashes, the exact declared
matrix and process order, execution sidecars, and all-output numerical replay.
Independent references use direct DFTs at small lengths and a separately
structured binary64 FFT at larger lengths; module replay also models DC
conditioning and display conversion. Ordinary magnitude-vector relative
$\ell_2$ and $\ell_\infty$ errors must be at most $3\times10^{-4}$ for float
or $10^{-10}$ for double; a zero reference requires zero output. These are
engineering acceptance budgets, not guarantees for each weak bin.

The long complete-module decay fixture separately flags the region affected
by Rack's flush-to-zero behavior and Fourier's display mapping. It uses the
predeclared \code{module-decay-ftz-v1} absolute floor, retains every vector,
and excludes flagged vectors from relative-accuracy claims. The floor is an
engineering allowance, not a proven error bound. This distinction matters:
relative error can reach one in that tail while the absolute error remains
below its recorded floor. Native allocator storage and allocation calls are
opaque to the separate C++ allocation probe. Appendix~\ref{app:reproducibility}
identifies the raw evidence and deterministic publication tools.
